\chapter{Manipulative Strategies and How They Are Caught}\label{ch:s2:manipulative-strategies-and-how-they-are-caught}

A trader placed large orders he never meant to fill, cancelled them within a fraction of a second, and traded the other side of the price move they
caused. In United States v. Coscia the Seventh Circuit upheld his conviction, the first under the Commodity Exchange Act's \omterm{def:s2:manipulative-strategies-and-how-they-are-caught:spoofing}{anti-spoofing} provision.
The evidence was in the order log: 24,814 large orders in three months of 2011, of which 0.5\% traded, against small orders on the other side of which
about 52\% were filled. This chapter describes each manipulative practice only as far as the public record does: how it works, why it is illegal or
contested, the case that established it, and the detector that finds it. On synthetic account-days, a detector that combines the two features of the
Coscia record catches 77\% of planted \omterm{def:s2:manipulative-strategies-and-how-they-are-caught:spoofing}{spoofing} at a 1\% false-positive rate. The order-to-trade ratio, often quoted as a \omterm{def:s2:manipulative-strategies-and-how-they-are-caught:spoofing}{spoofing} signal, catches none:
market makers cancel more than spoofers do. The build is \texttt{firm.surveil}.

\section{Spoofing and layering}

\begin{definition}[Spoofing]\label{def:s2:manipulative-strategies-and-how-they-are-caught:spoofing}
\emph{Spoofing}\index{spoofing} is bidding or offering with the intent to cancel the bid or offer before execution, so as to create a false
impression of supply or demand and trade against the price move it causes.
\end{definition}

\begin{definition}[Layering]\label{def:s2:manipulative-strategies-and-how-they-are-caught:layering}
\emph{Layering}\index{layering} is \omterm{def:s2:manipulative-strategies-and-how-they-are-caught:spoofing}{spoofing} with several orders placed at successive price levels on one side of the book, kept away from the best
price so that they are seen but unlikely to trade.
\end{definition}

\textbf{How it works.} In the Coscia opinion's account, a program placed a small order on one side and large orders on the other. The large orders
suggested supply or demand that did not exist, and moved the price towards the small order, which filled. The large orders were then cancelled,
and the process ran in reverse. Each cycle took about two-thirds of a second. It was repeated to more than 450,000 large orders and earned \$1.4
million. The CFTC described Sarao's version in the E-mini S\&P 500 future: from 2009 a modified trading platform layered four to six exceptionally
large sell orders, one price level apart and at least three or four levels from the best ask, and most were cancelled without trading.

\textbf{Why it is illegal.} The Dodd-Frank Act of 2010 added to the Commodity Exchange Act a prohibition on \omterm{def:s2:manipulative-strategies-and-how-they-are-caught:spoofing}{spoofing}, defined as bidding or offering with the
intent to cancel before execution. The EU's Market Abuse Regulation treats orders that give false or misleading signals as manipulation, including
orders placed and cancelled by algorithms, and lists as an indicator orders that change the best bid or offer or the book's representation and are
removed before they execute. The offence is the intent, not the cancellation: most orders are cancelled, and Coscia argued that high-frequency traders
cancel 98\% of theirs.

\textbf{The enforcement cases.} Coscia was convicted on all counts, sentenced to 36 months and lost his appeal in August 2017. The court found intent in
the program's design: it cancelled the large orders after a set time, if the small orders filled, or if a single large order filled. Sarao was charged
by the CFTC and the Justice Department in April 2015. The CFTC alleged that he used the \omterm{def:s2:manipulative-strategies-and-how-they-are-caught:layering}{layering} program on more than 400 trading days, made more than
\$40 million, and applied close to \$200 million of persistent selling pressure for two hours before the Flash Crash of 6 May 2010. In November 2016 he
consented to a penalty of \$25.7 million and disgorgement of \$12.9 million.

\begin{dated}{2026-09}{The law}\label{dat:s2:manipulative-strategies-and-how-they-are-caught:law}
In the United States, 7~U.S.C.~6c(a)(5)(C) makes it unlawful to engage in trading that is, is of the character of, or is commonly known to the trade as
\omterm{def:s2:manipulative-strategies-and-how-they-are-caught:spoofing}{spoofing}, defined as bidding or offering with the intent to cancel the bid or offer before execution. In the EU, Article~12 of the Market Abuse Regulation
defines market manipulation to include false or misleading signals, trading at the open or close that misleads investors, algorithmic orders that
initiate or exacerbate a trend, and the manipulation of benchmarks. Its Annex~I lists indicators, among them transactions with no change of
beneficial ownership and orders removed before execution after changing the book's appearance.
\end{dated}

\section{Momentum ignition and marking the close}

\begin{definition}[Momentum ignition]\label{def:s2:manipulative-strategies-and-how-they-are-caught:ignition}
\emph{Momentum ignition}\index{momentum ignition} is entering orders or trades intended to start or exaggerate a price trend so that other traders'
responses move the price further, and trading against the move they cause.
\end{definition}

\begin{definition}[Marking the close]\label{def:s2:manipulative-strategies-and-how-they-are-caught:close}
\emph{Marking the close}\index{marking the close} (banging the close) is trading heavily just before and during the closing period to move the
closing or settlement price, to benefit a position whose value depends on that price.
\end{definition}

\textbf{How they work.} The Market Abuse Regulation names orders that initiate or exacerbate a trend as manipulation; the chapter found no enforcement
case that uses the term \omterm{def:s2:manipulative-strategies-and-how-they-are-caught:ignition}{momentum ignition}, and the CFTC's Sarao complaint describes \omterm{def:s2:manipulative-strategies-and-how-they-are-caught:spoofing}{spoofing} designed to cause price swings that could be exploited.
\omterm{def:s2:manipulative-strategies-and-how-they-are-caught:close}{Marking the close} is better documented. In the CFTC's Optiver case, traders accumulated large Trading at Settlement positions, priced at the day's
settlement, in NYMEX crude oil, heating oil and gasoline futures. They then traded futures in the opposite direction shortly before and during the
close, to push the settlement in their favour.

\textbf{Why it is illegal.} Trading at the close that misleads investors relying on closing prices is manipulation under the Market Abuse Regulation,
and trading timed to move settlement and reference prices is one of its indicators. In the United States it is manipulation or attempted manipulation
under the Commodity Exchange Act.

\textbf{The enforcement case.} The complaint alleged 19 attempts in March 2007, at least 5 of them successful, and false statements to NYMEX. NYMEX's
own surveillance detected the trading. In April 2012 a consent order imposed a \$13 million civil penalty and \$1 million in disgorgement, trading
limits on Optiver, and trading bans of two to eight years on three former officers.

\section{Wash trades and benchmark manipulation}

\begin{definition}[Benchmark manipulation]\label{def:s2:manipulative-strategies-and-how-they-are-caught:benchmark}
\emph{Benchmark manipulation}\index{benchmark manipulation} is submitting false inputs to a benchmark, or trading or coordinating to move the data
it is calculated from, to benefit positions valued at the benchmark.
\end{definition}

\textbf{How they work.} A wash trade (Book~3, chapter~16) has the same beneficial owner on both sides, so it creates volume and prices without a change
of ownership. In October 2024 the SEC charged three firms that described themselves as market makers, and nine individuals, with selling market
manipulation as a service to crypto-asset promoters. They were alleged to have self-traded and used bots that at times generated billions of dollars of
artificial volume a day. Cong and co-authors estimated wash trading at over 70\% of reported volume on unregulated crypto exchanges. \omterm{def:s2:manipulative-strategies-and-how-they-are-caught:benchmark}{Benchmark
manipulation} is best known from the foreign exchange fixes. In May 2015 four banks agreed to plead guilty to conspiring to manipulate the euro-dollar
rate: their traders, in a chat room called ``The Cartel'', coordinated around the 1:15 p.m.\ ECB and 4:00 p.m.\ WM/Reuters fixes (Book~2,
chapter~17).

\textbf{Why they are illegal.} Transactions with no change of beneficial ownership are an indicator of manipulation under the Market Abuse Regulation,
and manipulating a benchmark's calculation is manipulation under its Article~12. In the United States, wash trades in futures are prohibited by the
Commodity Exchange Act, and securities manipulation by the securities laws.

\textbf{The enforcement cases.} The foreign exchange guilty pleas carried fines of more than \$2.5 billion. The SEC's 2024 case was part of an
investigation in which the FBI created a token of its own and the firms were alleged to have manipulated its market.

\section{Sandwiching and its legal status}

A sandwich attack (Book~3, chapter~22) buys ahead of a victim's pending swap on a decentralised exchange and sells right after it, in the same block,
so the victim trades at a worse price. It uses information that the blockchain makes public by design: pending transactions in a public mempool. Users
can avoid it by sending transactions through private channels, and services such as Flashbots Protect hide transactions from sandwich bots. Whether
sandwiching is front-running in the legal sense, manipulation or neither is unsettled: the chapter found no court ruling on it in the record it
searched, and it names no case.

\section{Surveillance: detectors and their errors}

\texttt{firm.surveil} builds account-days, one row per account per day, from legitimate types chosen to trip naive detectors, and plants 100 episodes
of each manipulation with the patterns described in the enforcement records. Among 11,500 legitimate account-days there are 5,000 market makers who cancel
almost everything, 3,000 who refresh all their quotes after every fill, 1,500 deep-book providers whose large orders far from the touch rarely fill,
and 2,000 directional traders. Each detector gets a threshold set so that 1\% of legitimate account-days are flagged
(\cref{lst:s2:manipulative-strategies-and-how-they-are-caught:detect}):

\begin{center}
\small
\begin{tabular}{@{}llcl@{}}
\toprule
practice & detector & caught & most flagged legitimate type\\
\midrule
\omterm{def:s2:manipulative-strategies-and-how-they-are-caught:spoofing}{spoofing} & order-to-trade ratio & 0\% & market makers (1.8\%)\\
 & fill-rate gap, small and large & 37\% & deep providers (7.7\%)\\
 & cancels after fills & 51\% & quote refreshers (3.8\%)\\
 & gap $\times$ cancels & 77\% & deep providers (7.7\%)\\
\omterm{def:s2:manipulative-strategies-and-how-they-are-caught:close}{marking the close} & share of volume at the close & 1\% & index funds (3.7\%)\\
 & share $\times$ next-day reversal & 21\% & index funds (3.7\%)\\
 & $\times$ settlement position & 61\% & index funds (3.3\%)\\
wash trades & self-matched trades & 74\% & multi-algorithm firms (4.0\%)\\
 & share of trades self-matched & 97\% & multi-algorithm firms (5.0\%)\\
\bottomrule
\end{tabular}
\end{center}

The order-to-trade ratio fails because the spoofer's defining behaviour is not cancelling a lot: market makers cancel more. It is cancelling one kind
of order and trading another, which is the Coscia testimony: two order sizes in use, with very different cancellation rates. Each half of that pattern
has innocent explanations. Deep-book providers' large orders rarely fill; quote refreshers cancel right after fills. Only together do the two features
separate the spoofers (\cref{fig:s2:manipulative-strategies-and-how-they-are-caught:roc}). Closing-window volume alone flags index funds, which must
trade at the close; reversal and a position priced at the settlement are what the Optiver record adds. Self-matches flag firms whose independent
algorithms occasionally cross, so the share of an account's trades that self-match works better than the count.

\begin{omfigure}
\begin{tikzpicture}
\begin{axis}[width=11cm,height=5.4cm,xmode=log,xlabel={\small false-positive rate (\%, log scale)},ylabel={\small spoofing episodes caught (\%)},
  tick label style={font=\footnotesize},xmin=0.1,xmax=50,ymin=0,ymax=100,xtick={0.1,1,10},xticklabels={0.1,1,10},grid=major,
  grid style={gray!25},table/col sep=comma,legend style={font=\scriptsize,at={(0.5,-0.3)},anchor=north},legend columns=2]
\addplot[omThm,very thick] table[x=fpr,y=both]{figdata/strategies-2/29-manipulative-strategies-and-how-they-are-caught/roc.csv};
\addplot[omDef,thick,dashed] table[x=fpr,y=cancel]{figdata/strategies-2/29-manipulative-strategies-and-how-they-are-caught/roc.csv};
\addplot[gray,thick] table[x=fpr,y=gap]{figdata/strategies-2/29-manipulative-strategies-and-how-they-are-caught/roc.csv};
\addplot[black,thick,dotted] table[x=fpr,y=otr]{figdata/strategies-2/29-manipulative-strategies-and-how-they-are-caught/roc.csv};
\legend{fill-rate gap $\times$ cancels after fills,cancels after fills,fill-rate gap,order-to-trade ratio}
\end{axis}
\end{tikzpicture}

\omcaption{Share of planted \omterm{def:s2:manipulative-strategies-and-how-they-are-caught:spoofing}{spoofing} account-days caught by four detectors against the share of legitimate account-days flagged. Data:
\texttt{s2\_surveil.roc}.}\label{fig:s2:manipulative-strategies-and-how-they-are-caught:roc}
\end{omfigure}

A surveillance alert is where an investigation starts, not the verdict. At a 1\% false-positive rate, 115 legitimate account-days are flagged for
every 77 true \omterm{def:s2:manipulative-strategies-and-how-they-are-caught:spoofing}{spoofing} episodes caught in this sample. Real surveillance ranks alerts, reads the orders behind them and looks for intent, as the
courts did (\cref{fig:s2:manipulative-strategies-and-how-they-are-caught:flagged}).

\begin{omfigure}
\begin{tikzpicture}
\begin{axis}[width=11cm,height=5.2cm,ybar,bar width=5pt,ylabel={\small flagged at 1\% false positives (\%)},tick label style={font=\footnotesize},
  xtick=data,xticklabels={market maker,refresher,deep provider,directional,spoofer},ymin=0,ymax=100,grid=major,grid style={gray!25},
  table/col sep=comma,legend style={font=\scriptsize,at={(0.03,0.97)},anchor=north west}]
\addplot[fill=gray!40,draw=gray] table[x=k,y=otr]{figdata/strategies-2/29-manipulative-strategies-and-how-they-are-caught/flagged.csv};
\addplot[fill=white,draw=black] table[x=k,y=gap]{figdata/strategies-2/29-manipulative-strategies-and-how-they-are-caught/flagged.csv};
\addplot[fill=omDef!50,draw=omDef] table[x=k,y=cancel]{figdata/strategies-2/29-manipulative-strategies-and-how-they-are-caught/flagged.csv};
\addplot[fill=omThm!60,draw=omThm] table[x=k,y=both]{figdata/strategies-2/29-manipulative-strategies-and-how-they-are-caught/flagged.csv};
\legend{order-to-trade,fill-rate gap,cancels after fills,both}
\end{axis}
\end{tikzpicture}

\omcaption{Share of each synthetic account type flagged by the four \omterm{def:s2:manipulative-strategies-and-how-they-are-caught:spoofing}{spoofing} detectors at thresholds that flag 1\% of all legitimate account-days.
Data: \texttt{s2\_surveil.detectors}.}\label{fig:s2:manipulative-strategies-and-how-they-are-caught:flagged}
\end{omfigure}

\section{Strategy files}

The seven manipulative practices below are described, in the order this series requires, by how they work, why they are illegal or contested, the
enforcement case and the surveillance detector. The last file is the surveillance desk's own.

\begin{strategyfile}{Spoofing}\label{strat:s2:manipulative-strategies-and-how-they-are-caught:spoofing}
\sfield{Who is harmed} Traders who respond to displayed orders and trade at prices moved by orders that were never meant to trade.
\sfield{Instruments and venues} Futures and other markets with visible central limit order books.
\sfield{How it works} Large orders on one side, cancelled before they can trade, and small orders on the other that fill after the price moves.
\sfield{Why it is illegal} Bidding or offering with intent to cancel before execution: 7~U.S.C.~6c(a)(5)(C); false or misleading signals under the
Market Abuse Regulation.
\sfield{Enforcement case} United States v. Coscia (7th Cir. 2017): 36 months, conviction affirmed.
\sfield{Surveillance detector} Small and large orders' fill rates compared, with large cancellations right after fills on the other side.
\sfield{How it is caught} Exchange and regulator surveillance of order logs; intent read from program design and trading records.
\sfield{Evaluating the detector honestly} Legitimate look-alikes in the evaluation (deep-book providers, quote refreshers); false positives at the
threshold used.
\sfield{Sources} Seventh Circuit opinion (2017); this chapter: 77\% caught at 1\% false positives.
\end{strategyfile}

\begin{strategyfile}{Layering}\label{strat:s2:manipulative-strategies-and-how-they-are-caught:layering}
\sfield{Who is harmed} Traders reading the depth of the book.
\sfield{Instruments and venues} Deep, visible order books; the E-mini S\&P 500 in the Sarao case.
\sfield{How it works} Several large orders at successive levels, kept away from the best price and cancelled, as in the CFTC's account.
\sfield{Why it is illegal} \omterm{def:s2:manipulative-strategies-and-how-they-are-caught:spoofing}{Spoofing} under US law; orders that change the book's representation and are removed before execution under the Market Abuse
Regulation.
\sfield{Enforcement case} CFTC v. Nav Sarao Futures Limited PLC and Sarao (2015): penalty and disgorgement of \$38.6 million by consent (2016).
\sfield{Surveillance detector} Persistent large orders several levels from the touch, re-priced as the market moves, with a near-zero fill rate.
\sfield{How it is caught} Order-book reconstruction; complaints from other participants.
\sfield{Evaluating the detector honestly} Market makers who quote size deep in the book look the same until trading on the other side is added.
\sfield{Sources} CFTC press releases 7156-15 and 7486-16.
\end{strategyfile}

\begin{strategyfile}{Momentum ignition}\label{strat:s2:manipulative-strategies-and-how-they-are-caught:momentum}
\sfield{Who is harmed} Traders whose algorithms follow short-term trends.
\sfield{Instruments and venues} Electronic markets with trend-following participants.
\sfield{How it works} Orders or trades intended to start or exaggerate a move, traded against as others follow.
\sfield{Why it is illegal} Market Abuse Regulation Art. 12(2)(c)(iii): orders that initiate or exacerbate a trend.
\sfield{Enforcement case} No case using the term was found; \omterm{def:s2:manipulative-strategies-and-how-they-are-caught:spoofing}{spoofing} cases describe price swings caused to be exploited.
\sfield{Surveillance detector} Bursts of aggressive orders followed by opposite trading and a reversal (Annex~I indicator (e)).
\sfield{How it is caught} Pattern alerts on concentrated trading with reversals.
\sfield{Evaluating the detector honestly} Legitimate execution algorithms also trade in bursts; reversals are common.
\sfield{Sources} Market Abuse Regulation, Art. 12 and Annex~I.
\end{strategyfile}

\begin{strategyfile}{Wash trades}\label{strat:s2:manipulative-strategies-and-how-they-are-caught:wash}
\sfield{Who is harmed} Investors who read volume and prices as genuine interest.
\sfield{Instruments and venues} Crypto assets on unregulated venues; futures; thinly traded securities.
\sfield{How it works} Trading with oneself or a related account, with no change of beneficial ownership.
\sfield{Why it is illegal} Prohibited in US futures and securities markets; an indicator of manipulation under the Market Abuse Regulation.
\sfield{Enforcement case} SEC charges against three purported market makers and nine individuals (October 2024).
\sfield{Surveillance detector} The share of an account's trades matched against the same beneficial owner.
\sfield{How it is caught} Beneficial-ownership data across accounts; exchange self-match prevention records.
\sfield{Evaluating the detector honestly} Firms with independent algorithms cross each other by accident; counts penalise large honest traders.
\sfield{Sources} SEC press release 2024-166; Cong and co-authors (2022); this chapter: 97\% caught at 1\% false positives.
\end{strategyfile}

\begin{strategyfile}{Marking the close}\label{strat:s2:manipulative-strategies-and-how-they-are-caught:close}
\sfield{Who is harmed} Holders of positions and contracts priced at the close or settlement.
\sfield{Instruments and venues} Futures settlements; closing auctions; Trading at Settlement contracts.
\sfield{How it works} A position priced at the settlement, then heavy trading against it in the closing minutes, as in the CFTC's Optiver account.
\sfield{Why it is illegal} Manipulation under the Commodity Exchange Act; trading at the close that misleads investors under the Market Abuse
Regulation.
\sfield{Enforcement case} CFTC v. Optiver (consent order 2012): \$14 million; trading bans for former officers.
\sfield{Surveillance detector} Closing-window share of volume, with a reversal the next day and a position priced at the settlement.
\sfield{How it is caught} Exchange surveillance of settlement windows; NYMEX detected Optiver's trading.
\sfield{Evaluating the detector honestly} Index funds must trade at the close; volume alone flags them.
\sfield{Sources} CFTC press release 6239-12; this chapter: 61\% caught at 1\% false positives.
\end{strategyfile}

\begin{strategyfile}{Benchmark manipulation}\label{strat:s2:manipulative-strategies-and-how-they-are-caught:benchmark}
\sfield{Who is harmed} Everyone whose contracts reference the benchmark.
\sfield{Instruments and venues} Foreign exchange fixes; interest-rate benchmarks; commodity assessments.
\sfield{How it works} False submissions, or coordinated trading around the calculation window.
\sfield{Why it is illegal} Market Abuse Regulation Art. 12(1)(d); price fixing and manipulation under US law.
\sfield{Enforcement case} Four banks' guilty pleas for foreign exchange fix manipulation (May 2015), fines over \$2.5 billion.
\sfield{Surveillance detector} Trading concentrated around calculation times, with communications review.
\sfield{How it is caught} Communications surveillance; benchmark administrators' data checks.
\sfield{Evaluating the detector honestly} Clients' fix orders are legitimately executed around the fix.
\sfield{Sources} US Department of Justice (2015); Book~2, chapter~17.
\end{strategyfile}

\begin{strategyfile}{Sandwiching}\label{strat:s2:manipulative-strategies-and-how-they-are-caught:sandwich}
\sfield{Who is harmed} Users whose pending swaps are visible in a public mempool.
\sfield{Instruments and venues} Decentralised exchanges on public blockchains.
\sfield{How it works} A purchase before a pending swap and a sale after it, in the same block (Book~3, chapter~22).
\sfield{Why it is contested} It exploits information public by design; no court ruling on it was found in the record searched.
\sfield{Enforcement case} None named.
\sfield{Surveillance detector} Pairs of transactions by one address bracketing another's swap in the same pool and block.
\sfield{How it is caught} It is avoided rather than caught: private transaction channels hide swaps.
\sfield{Evaluating the detector honestly} Arbitrage that happens to bracket a swap looks alike.
\sfield{Sources} Flashbots documentation.
\end{strategyfile}

\begin{strategyfile}{Detector: cancellation-pattern surveillance}\label{strat:s2:manipulative-strategies-and-how-they-are-caught:detector}
\sfield{Who pays you, and why} The venue, the broker or the firm itself, obliged to detect and report suspicious orders.
\sfield{Instruments and venues} Every order log the firm or venue keeps.
\sfield{Signal} Per account-day: the gap between small and large orders' fill rates, times the share of large cancellations right after a fill on the
other side.
\sfield{Sizing and execution} A threshold set to the number of alerts investigators can read.
\sfield{Costs} Investigators' time; false positives.
\sfield{How it dies} Manipulators who mix sizes and timings; legitimate strategies that drift into the pattern.
\sfield{Horizon, capacity, infrastructure} Daily; complete, time-stamped order logs.
\sfield{Backtest honestly} Evaluate on legitimate look-alikes, at the operating threshold, with the episodes labelled from real cases where possible.
\sfield{Sources} United States v. Coscia (2017); this chapter: 77\% at 1\% false positives, against 0\% for the order-to-trade ratio.
\end{strategyfile}

\section{Tutorial: in the order log}\label{tut:s2:manipulative-strategies-and-how-they-are-caught:tut}

\begin{tutorial}
\textbf{Goal.} Build account-days of legitimate types and planted episodes, score them with detectors drawn from the enforcement records, and measure
each detector's catch rate at a fixed false-positive rate. \textbf{End state:} the table and two figures.
\begin{tutsteps}
\item \textbf{The \omterm{def:s2:manipulative-strategies-and-how-they-are-caught:spoofing}{spoofing} detectors}.
\omcode{code/firm/surveil/firm_surveil.py}{67}{76}{Order-to-trade, fill-rate gap, cancels after fills, and both.}\label{lst:s2:manipulative-strategies-and-how-they-are-caught:detect}
\item \textbf{Evaluation}.
\omcode{code/firm/surveil/firm_surveil.py}{115}{124}{The catch rate at the threshold that flags a fixed share of legitimate accounts.}\label{lst:s2:manipulative-strategies-and-how-they-are-caught:eval}
\item \textbf{Run} \texttt{detectors()}, \texttt{roc()} and \texttt{fig\_surveil.py}.
\end{tutsteps}
\textbf{What to change next.} Add a legitimate type that hedges options with small orders against large quotes; measure how many alerts per day an
investigator would read; add the closing detector's position data from a second source.
\end{tutorial}

\section{Build: surveillance}\label{bld:s2:manipulative-strategies-and-how-they-are-caught:surveil}

\begin{build}
\bfield{Purpose} Labelled account-days and detectors for \omterm{def:s2:manipulative-strategies-and-how-they-are-caught:spoofing}{spoofing}, \omterm{def:s2:manipulative-strategies-and-how-they-are-caught:close}{marking the close} and wash trades, evaluated at fixed false-positive rates.
\bfield{Interface} \texttt{SurveilConfig(\dots)}, \texttt{spoofing\_days}, \texttt{close\_days}, \texttt{wash\_days}, \texttt{spoof\_scores},
\texttt{close\_scores}, \texttt{wash\_scores}, \texttt{tpr\_at\_fpr(score, label, fpr)}, \texttt{flagged\_by\_type}.
\bfield{Rules} Legitimate types chosen to trip each detector; episodes only as described in enforcement records; no manipulation's profit modelled.
\bfield{Acceptance tests} \texttt{code/firm/surveil/tests/}: the catch rate by hand; scores on two account-days; the combined detector beats its parts
and the order-to-trade ratio fails; the self-match share by hand.
\bfield{Stretch} Order-level logs; ranking alerts; detectors learned from labelled cases.
\end{build}

\begin{omsources}
\item United States v. Coscia, No.~16-3017 (7th Cir. 2017), 866 F.3d 782.
\item Commodity Futures Trading Commission, press releases 7156-15 (21 April 2015), 7486-16 (17 November 2016) and 6239-12 (19 April 2012).
\item Regulation (EU) No 596/2014 on market abuse, Article~12 and Annex~I.
\item US Department of Justice, press release on foreign exchange guilty pleas, 20 May 2015.
\item Securities and Exchange Commission, press release 2024-166, 9 October 2024.
\item L.~W.~Cong, X.~Li, K.~Tang and Y.~Yang, ``Crypto wash trading'', NBER Working Paper 30783, 2022.
\end{omsources}

\section{Exercises}

\begin{exercise}[$\star$]\label{exo:s2:manipulative-strategies-and-how-they-are-caught:1}
A trader placed 24,814 large orders, of which 0.5\% traded. How many traded?
\end{exercise}

\begin{exercise}[$\star$]\label{exo:s2:manipulative-strategies-and-how-they-are-caught:2}
At a 1\% false-positive rate on 11,500 legitimate account-days, how many are flagged, and how many true episodes of 100 does a detector with a 77\%
catch rate find?
\end{exercise}

\begin{exercise}[$\star$]\label{exo:s2:manipulative-strategies-and-how-they-are-caught:3}
Sarao's consent order required a penalty of \$25.7 million and disgorgement of \$12.9 million. What was the total?
\end{exercise}

\begin{exercise}[$\star\star$]\label{exo:s2:manipulative-strategies-and-how-they-are-caught:4}
Why does the order-to-trade ratio fail to catch spoofers?
\end{exercise}

\begin{exercise}[$\star\star$]\label{exo:s2:manipulative-strategies-and-how-they-are-caught:5}
Why is closing-window volume alone a poor detector of \omterm{def:s2:manipulative-strategies-and-how-they-are-caught:close}{marking the close}?
\end{exercise}

\begin{exercise}[$\star\star$]\label{exo:s2:manipulative-strategies-and-how-they-are-caught:6}
Why must intent, not cancellation, define \omterm{def:s2:manipulative-strategies-and-how-they-are-caught:spoofing}{spoofing}?
\end{exercise}

\begin{exercise}[$\star\star\star$]\label{exo:s2:manipulative-strategies-and-how-they-are-caught:7}
\emph{Coding.} Run \texttt{roc()}. How many \omterm{def:s2:manipulative-strategies-and-how-they-are-caught:spoofing}{spoofing} episodes does the combined detector catch at false-positive rates of 0.1\% and 10\%, and what
does that mean for an investigator's workload?
\end{exercise}

\begin{exercise}[$\star\star\star$]\label{exo:s2:manipulative-strategies-and-how-they-are-caught:8}
\emph{Find the flaw.} ``Our surveillance flagged no one with an order-to-trade ratio above our limit this year, so we have no \omterm{def:s2:manipulative-strategies-and-how-they-are-caught:spoofing}{spoofing}.''
\end{exercise}

\section{Problem: In the Order Log}

\begin{problem}[{Weekend problem --- manipulative strategies and surveillance}]\label{pb:s2:manipulative-strategies-and-how-they-are-caught:1}
The public enforcement record and the chapter's synthetic account-days.

\medskip\noindent\textbf{Part I --- \omterm{def:s2:manipulative-strategies-and-how-they-are-caught:spoofing}{Spoofing}.}
\begin{enumerate}
\item Define \omterm{def:s2:manipulative-strategies-and-how-they-are-caught:spoofing}{spoofing} and \omterm{def:s2:manipulative-strategies-and-how-they-are-caught:layering}{layering}.
\item What did the Coscia record show about his orders?
\item What did the CFTC allege about Sarao?
\item What does US law prohibit, and what does the Market Abuse Regulation list?
\end{enumerate}

\medskip\noindent\textbf{Part II --- Other practices.}
\begin{enumerate}[resume]
\item Define \omterm{def:s2:manipulative-strategies-and-how-they-are-caught:ignition}{momentum ignition} and \omterm{def:s2:manipulative-strategies-and-how-they-are-caught:close}{marking the close}.
\item What happened in the Optiver case?
\item Define \omterm{def:s2:manipulative-strategies-and-how-they-are-caught:benchmark}{benchmark manipulation}; what happened at the FX fixes?
\item What did the SEC allege in 2024?
\end{enumerate}

\medskip\noindent\textbf{Part III --- Surveillance.}
\begin{enumerate}[resume]
\item Describe the synthetic account types and why each is there.
\item Give the detector table.
\item Why does combining two features work?
\item What is the legal status of sandwiching?
\end{enumerate}

\medskip\noindent\textbf{Part IV --- The verdict.}
\begin{enumerate}[resume]
\item State the \emph{named result}: each detector's true-positive rate at a fixed false-positive rate, and how legitimate market making triggers them.
\item What happens to an alert after the detector fires?
\item How would you evaluate a detector honestly?
\item What would a manipulator change to avoid the combined detector, and what would that cost the manipulation?
\item Why do the enforcement cases matter to detector design?
\item Which strategy file is legitimate?
\item How does this chapter relate to chapter~24's mark-outs?
\item In one sentence: what separates a market maker from a spoofer?
\end{enumerate}
\end{problem}

\section{Interview questions}

\begin{interviewq}[$\star$ \iqroles{trader}]\label{iq:s2:manipulative-strategies-and-how-they-are-caught:1}
You cancel 95\% of your orders. How do you show you are not \omterm{def:s2:manipulative-strategies-and-how-they-are-caught:spoofing}{spoofing}?
\end{interviewq}

\begin{interviewq}[$\star\star$ \iqroles{researcher}]\label{iq:s2:manipulative-strategies-and-how-they-are-caught:2}
Design a \omterm{def:s2:manipulative-strategies-and-how-they-are-caught:spoofing}{spoofing} detector from an order log, and say how you would test it.
\end{interviewq}

\begin{interviewq}[$\star\star$ \iqroles{trader}]\label{iq:s2:manipulative-strategies-and-how-they-are-caught:3}
You hold a large position priced at today's settlement. What may you do in the closing minutes?
\end{interviewq}

\begin{interviewq}[$\star\star$ \iqroles{risk}]\label{iq:s2:manipulative-strategies-and-how-they-are-caught:4}
What controls should a trading firm have against its own algorithms manipulating markets?
\end{interviewq}

\begin{interviewq}[$\star\star$ \iqroles{developer}]\label{iq:s2:manipulative-strategies-and-how-they-are-caught:5}
What must an order log record for surveillance to work?
\end{interviewq}

\begin{interviewq}[$\star\star\star$ \iqroles{researcher}]\label{iq:s2:manipulative-strategies-and-how-they-are-caught:6}
A detector flags 1\% of 11,500 legitimate account-days and 77 of 100 episodes. What share of alerts are true, and what happens to it if episodes
are ten times rarer?
\end{interviewq}
